\documentclass[tikz,a4paper,11pt]{book}
%\documentclass[a4paper,twoside,11pt,titlepage]{book}
\usepackage{listings}
\usepackage[utf8]{inputenc}
\usepackage{csquotes}
\usepackage[spanish,es-noquoting]{babel}
\usepackage{forest}
\usepackage{xcolor}
\usepackage[export]{adjustbox}
\usepackage{caption}
\usepackage{float}
\usetikzlibrary{arrows.meta,shapes,positioning,trees}

% \usepackage[style=list, number=none]{glossary} %
%\usepackage{titlesec}
%\usepackage{pailatino}

\decimalpoint
\usepackage{dcolumn}
\newcolumntype{.}{D{.}{\esperiod}{-1}}
\makeatletter
\addto\shorthandsspanish{\let\esperiod\es@period@code}
\makeatother


%\usepackage[chapter]{algorithm}
\RequirePackage{verbatim}
%\RequirePackage[Glenn]{fncychap}
\usepackage{fancyhdr}
\usepackage{graphicx}
\usepackage{afterpage}
\usepackage{longtable}
\usepackage[pdfborder={000}]{hyperref} %referencia


% ********************************************************************
% Re-usable information
% ********************************************************************
\newcommand{\myTitle}{Chief\xspace}
\newcommand{\myDegree}{Grado en Ingeniería Informática\xspace}
\newcommand{\myName}{Sergio Cervilla Ortega\xspace}
\newcommand{\myProf}{Juan Julián Merelo Guervós\xspace}
\newcommand{\myFaculty}{Escuela Técnica Superior de Ingenierías Informática y de Telecomunicación\xspace}
\newcommand{\myFacultyShort}{E.T.S. de Ingenierías Informática y de Telecomunicación\xspace}
% \newcommand{\myDepartment}{Departamento de Arquitectura y Tecnología de Computadores\xspace}
\newcommand{\myUni}{\protect{Universidad de Granada}\xspace}
\newcommand{\myLocation}{Granada\xspace}
\newcommand{\myTime}{\today\xspace}
\newcommand{\myVersion}{Version 1.0\xspace}
\newcommand{\innerxsep}{5mm}

\hypersetup{
pdfauthor = {\myName (cervick13@gmail.com)},
pdftitle = {\myTitle},
pdfsubject = {},
pdfcreator = {LaTeX con el paquete ....},
pdfproducer = {pdflatex}
}

%\hyphenation{}


%\usepackage{doxygen/doxygen}
%\usepackage{pdfpages}
\usepackage{url}
\usepackage{hyperref}
\usepackage{colortbl,longtable}
\usepackage[stable]{footmisc}
%\usepackage{index}

%\makeindex
%\usepackage[style=long, cols=2,border=plain,toc=true,number=none]{glossary}
% \makeglossary

% Definición de comandos que me son tiles:
%\renewcommand{\indexname}{Índice alfabético}
%\renewcommand{\glossaryname}{Glosario}

\pagestyle{fancy}
\fancyhf{}
\fancyhead[LO]{\leftmark}
\fancyhead[RE]{\rightmark}
\fancyhead[RO,LE]{\textbf{\thepage}}


%Para conseguir que en las páginas en blanco no ponga cabecerass
\makeatletter
\def\clearpage{%
  \ifvmode
    \ifnum \@dbltopnum =\m@ne
      \ifdim \pagetotal <\topskip
        \hbox{}
      \fi
    \fi
  \fi
  \newpage
  \thispagestyle{empty}
  \write\m@ne{}
  \vbox{}
  \penalty -\@Mi
}
\makeatother


\usepackage{pdfpages}
\usepackage{fourier}
\begin{document}
\begin{titlepage}
\newlength{\centeroffset}
\setlength{\centeroffset}{-0.5\oddsidemargin}
\addtolength{\centeroffset}{0.5\evensidemargin}
\thispagestyle{empty}

\noindent\hspace*{\centeroffset}\begin{minipage}{\textwidth}

\centering
\includegraphics[width=0.9\textwidth]{imagenes/logo_ugr.jpg}\\[1.4cm]

\textsc{ \Large TRABAJO FIN DE GRADO\\[0.2cm]}
\textsc{ GRADO EN INGENIERIA INFORMATICA}\\[1cm]

{\Huge\bfseries Chief \\}
\noindent\rule[-1ex]{\textwidth}{3pt}\\[3.5ex]
{\large\bfseries Apoyo a la labor de la Policía Local }
\end{minipage}

\vspace{2.5cm}
\noindent\hspace*{\centeroffset}
\begin{minipage}{\textwidth}
\centering

\textbf{Autor}\\ {Sergio Cervilla Ortega}\\[2.5ex]
\textbf{Director}\\ {Juan Julián Merelo Guervós}\\[2cm]
\includegraphics[width=0.3\textwidth]{imagenes/etsiit_logo.png}\\[0.1cm]
\textsc{Escuela Técnica Superior de Ingenierías Informática y de Telecomunicación}\\
\textsc{---}\\
Granada, algún mes de 2018
\end{minipage}
\end{titlepage}

\chapter*{}
\thispagestyle{empty}

\begin{center}
{\large\bfseries Chief \\ Apoyo a la labor de la Policía Local}\\
\end{center}
\begin{center}
Sergio Cervilla Ortega\\
\end{center}

%\vspace{0.7cm}

\noindent{\textbf{Resumen}}\\

Lorem ipsum dolor sit amet, consectetur adipisicing elit, sed do eiusmod tempor incididunt ut labore et dolore magna aliqua. Ut enim ad minim veniam, quis nostrud exercitation ullamco laboris nisi ut aliquip ex ea commodo consequat. Duis aute irure dolor in reprehenderit in voluptate velit esse cillum dolore eu fugiat nulla pariatur. Excepteur sint occaecat cupidatat non proident, sunt in culpa qui officia deserunt mollit anim id est laborum.

\cleardoublepage

\thispagestyle{empty}

\noindent\rule[-1ex]{\textwidth}{2pt}\\[4.5ex]

D. \textbf{Juan Julián Merelo Guervós}, Profesor del Área de XXXX del Departamento de Arquitectura y Tecnología de Computadores de la Universidad de Granada.

\vspace{0.5cm}

\textbf{Informo:}

\vspace{0.5cm}

Que el presente trabajo, titulado \textit{\textbf{Chief, blah blah}},
ha sido realizado bajo mi supervisión por \textbf{Sergio Cervilla Ortega}, y autorizo la defensa de dicho trabajo ante el tribunal
que corresponda.

\vspace{0.5cm}

Y para que conste, expiden y firman el presente informe en Granada a X de mes de 2018.

\vspace{1cm}

\textbf{El director: }

\vspace{5cm}

\noindent \textbf{Juan Julián Merelo Guervós}

\frontmatter
\tableofcontents
\listoffigures
%\listoftables
%\mainmatter
%\setlength{\parskip}{5pt}

\chapter{Introducción}

Vivimos cada día en un mundo más conectado a Internet, es una época en la que se está dejando de 
lado el lápiz y el papel. Donde se está abriendo paso la tecnología. Estos nuevos medios, bien
usados, son capaces de facilitarnos en numerosas ocasiones las tareas del día a día. Convierte acciones
tediosas en simples gestos que, con el tiempo, acabamos automatizando y realizamos sin esfuerzo. 
Bajo dicha motivación, decidí crear un sistema informatizado para ayudar en las labores del cuerpo policial de Maracena.\\

Previamente, el cuerpo policial utilizaba un esquema bastante anticuado para desemplear sus funciones.
Por ejemplo, los registros de incidencias se iban apuntando manualmente y al final del turno
se reunían para poner en común todos los datos en un documento que posteriormente se guardaba.
Y al igual que en el ejemplo anterior, debían esperar a llegar a las dependencias para poder 
rellenar cualquier tipo de denuncia administrativa con los datos que habían apuntado en 
un papel para no olvidarlos.\\

Para ello, he diseñado una aplicación web que se encarga de gestionar una gran cantidad de 
modelos de denuncia y que permite que se rellenen de forma segura, a prueba de 
fallos y de una manera muy intuitiva para el usuario. Además de poder gestionar dichos documentos,
se ha creado un sistema de registro de incidencias completamente automatizado para que en 
ningún momento se pierda ningún dato y un gran número de funcionalidades más que se explicarán en 
posteriores puntos. El Ayuntamiento de Maracena será pionero en este ámbito, ya que han sido los 
primeros que han apostado por las nuevas tecnologías. Una gran apuesta, pero que traerá grandes
beneficios a la calidad del trabajo del cuerpo policial y a la seguridad de los ciudadanos. \\

Este Trabajo de Fin de Grado ha sido realizado completamente por mi, Sergio 
Cervilla Ortega durante el año 2018 y es libre bajo la licencia GNU General Public License v3.0 \cite{gplv3}\\ 

Ha sido desarrollado bajo el permiso, aceptación, revisión y ayuda del Ayuntamiento de Maracena, localidad en
la que resido a fecha de hoy.\\


\chapter{Elección de la plataforma de gestión documental}

Esta es una de las elecciones más significativas debido a que toda la gestión documental
y la manera en la que se realiza depende del software que utilicemos.
\\

En este caso en particular vamos a utilizar \textbf{Nuxeo} debido a que se adapta mejor a mis
necesidades que mi segunda opción, \textbf{Alfresco}.
\\

Los motivos por los que he elegido Nuxeo frente a Alfresco son los siguientes:

\begin{itemize}
	\item Alfresco sigue la filosofía \textit{Open Source} como método de trabajo, sino que la utilizan a modo de marketing
	 			ya que podemos comprobar que en su desarrollo solo intervienen trabajadores de dicha empresa a pesar de ser bastante
				más grande que Nuxeo.
	\item Alfresco tiene dos versiones, \textit{Community edition} y \textit{Commercial edition}. Esto no sería un problema
				si en la versión de la comunidad se incluyera \textbf{toda la funcionalidad}. Pero de esta manera, si quieres tener
				el software al completo deberás pagar la versión comercial. Sin embargo, Nuxeo nos proporciona toda la funcionalidad
				base tengas o no tengas contratada una suscripción.
	\item La tecnología que impulsa Nuxeo es también de código libre y actual.
	\item La \textit{API} de Nuxeo es mucho más completa.
	\item La documentación de Nuxeo es muy completa, fácil de acceder y de aplicar.
\end{itemize}

% \chapter{Estado del arte}

Actualmente el estado del mercado para este tipo de desarrollos está muy limitado dado que 
la mayoría de los datos que manejan dichos sistemas son de carácter sensible. Por tanto, se 
debe tener un especial cuidado con el tratamiento de los mismos tanto a la hora de transmitirlos
como en su almacenamiento.\\

Esta premisa ha conseguido que la única aplicación actual del mercado, y aún en fase de desarrollo,
utilizada por la Guardia Civil de Tráfico sea de código privativo y una caja negra para el usuario.
De hecho, está en una fase tan temprana de su desarrollo que aún carece de nombre y de prácticamente 
utilidad.\\ 

insertar más datos de la GC

% [hablar de fuera de españa]

% [hablar del open source]
% \input{secciones/servicio_mensajeria}
% \input{secciones/planificacion}
% \input{secciones/implementacion}
% \input{secciones/conclusiones}

%
\nocite{*}
\bibliography{bibliografia}\addcontentsline{toc}{chapter}{Bibliografía}
\bibliographystyle{unsrt}


\end{document}
